\chapter{Usage of the package}
In this section, we provide a few minimal \python
scripts demonstrating the use of the \gretlpy module.
More comprehensive examples can be found in the
\texttt{examples} subdirectory of each \gretlpy distribution:
\begin{description}
	\item[\texttt{data\_management}:] Examples for managing
		datasets (including helper functions): \texttt{df2dict.py},
		\texttt{longley.csv}, \texttt{nparray2dict.py},
		\texttt{sample.py}, and \texttt{series2dict.py}.
	\item[\texttt{estimators}:] Examples of native estimators,
		as outlined in~\ref{sec:estimators_python}: \texttt{ar.py},
		\texttt{ar1.py}, \texttt{arima.py}, \texttt{biprobit.py},
		\texttt{duration.py}, \texttt{garch.py}, \texttt{heckit.py},
		\texttt{lad.py}, \texttt{logistic.py}, \texttt{logit.py},
		\texttt{mpols.py}, \texttt{negbin.py}, \texttt{ols.py},
		\texttt{panel.py}, \texttt{poisson.py}, \texttt{probit.py},
		\texttt{quantreg.py}, \texttt{tobit.py}, \texttt{tsls.py},
		\texttt{var.py}, \texttt{vecm.py}, and \texttt{wls.py}.
	\item[\texttt{module\_related\_tools}:] Examples for managing
		\gretlpy itself: \texttt{set.py}.
	\item[\texttt{statsmodels}:] Examples comparing \gretlpy with
		the \pythcode{statsmodels} package (see \ref{sec:examples.statsmodels}
		for details).
	\item[\texttt{Principles\_of\_Econometrics\_4th\_edition}:] All
		examples used in the great handbook \textit{Principles of Econometrics,
		Fourth Edition} (see \url{https://principlesofeconometrics.com/poe4/poe4.htm}
		for details): \texttt{chapter\_5.py}, \texttt{chapter\_6.py},
		\texttt{chapter\_7.py}, \texttt{chapter\_8.py}, \texttt{chapter\_9.py},
		\texttt{chapter\_10.py}, \texttt{chapter\_11.py}, \texttt{chapter\_12.py},
		\texttt{chapter\_13.py}, \texttt{chapter\_14.py}, \texttt{chapter\_15.py},
		and \texttt{chapter\_16.py},
		\begin{description}
			\item[\texttt{data}:] Subdirectory containing datasets
				in \gretl \texttt{gdt} format (huge thanks to
				Lee Addkins, \url{https://learneconometrics.com/}).
			\item[\texttt{functions}:] The \texttt{hausman\_taylor.gfn},
				a \gretl package by Allin Cottrell which implements
				the Hausman-Taylor estimator.
		\end{description}
\end{description}

All files mentioned above (excluding those in the \texttt{statsmodels}
subdirectory) define a \pythcode{run\_example()} function, whose body
contains the complete example script. Thus, to run the full example,
the user simply executes this function. Below are two useful snippets
demonstrating how to print the contents of a script and how to run it.
In both cases, the \pythcode{inspect} package is required.

To print the contents of the \texttt{ar.py} file, run:
\begin{lstlisting}[style=pythonstylefoot]
import inspect
import gretl.examples.estimators.ar

print(inspect.getsource(gretl.examples.estimators.ar.run_example))
\end{lstlisting}
The contents of the \texttt{ar.py} script are shown in
Figure~\ref{fig:estimators.ar.contents}:

\begin{figure}[h]
	\centering
	\includegraphics[width=0.8\columnwidth,keepaspectratio]{estimators.ar.contents}
	\caption{Contents of the \texttt{ar.py} script.}
	\label{fig:estimators.ar.contents}
\end{figure}

To execute the examples in the \texttt{ar.py} file, run:
\begin{lstlisting}[style=pythonstylefoot]
import inspect
import gretl.examples.estimators.ar

gretl.examples.estimators.ar.run_example()
\end{lstlisting}

Another set of examples, including those related to features
under development, can be found in the source tree of the
\gretlpy package (see
\url{https://sourceforge.net/p/gretl/gretl4py/ci/master/tree/demo/}).
Although the contents of the \texttt{demo} directory are
volatile by definition, some scripts are expected to remain
there indefinitely:

\begin{description}
	\item[\texttt{data\_management.py}:] A script containing
		examples of functions for dataset management (see
		Sections~\ref{sec:classpydataset} and
		\ref{sec:data_management_python}).
	\item[\texttt{estimators.py}:] A script demonstrating
		the use of various estimators provided by \gretlpy
		(see Sections~\ref{sec:estimators_python} and
		\ref{sec:estimators_python_examples}), including
		minimal results verification.
	\item[\texttt{module\_related\_tools.py}:] A script
		showcasing functions for configuring specific \gretlpy
		internal options (see Sections~\ref{sec:module_related_python}
		and \ref{sec:module_related_python_examples}).
	\item[\texttt{speed}:] A directory containing scripts
		to compare the performance of \gretlpy with other
		\python packages:
		\begin{description}
			\item[\texttt{arma\_speed.py}:] A script for
				comparing the performance of $ARIMA(4,0,4)$
				models. Typically, \gretlpy is approximately
				\textbf{50 times} faster than \pkg{statsmodels}.\newline
				\textbf{Note:} Before running the test, the dataset
				file must be generated by executing
				\texttt{gretlcli(.exe) -b genarma22\_test.inp}.
			\item[\texttt{probit\_speed.py}:] A script for testing
				the performance of a probit model with 24 regressors.
				In general, \gretlpy is about 20\% faster than
				\pkg{statsmodels}.\newline
				\textbf{Note:} Before running the test, the dataset
				file must be generated by executing
				\texttt{gretlcli(.exe) -b genprobit\_test.inp}.
	\end{description}
\end{description}
